\section{Impact Statement}
% The primary objective of this paper is to inform the visual generation community that the distributional matching between feature vectors and code vectors represents the optimal solution for vector quantization (VQ). Although the proposed quadratic Wasserstein distance may not be the most ideal distribution matching approach, we encourage researchers in this field to investigate and identify more effective distributional matching methods to ultimately achieve truly optimal VQ strategies. Given the technical nature of our work, we do not anticipate any significant societal consequences that require specific emphasis in this statement.

% The primary objective of this paper is to advance the understanding of vector quantization (VQ) by demonstrating that distributional matching between feature vectors and code vectors represents an optimal solution. This insight can benefit the broader visual generation community, including applications in image synthesis, data compression, and generative modeling. Although the proposed quadratic Wasserstein distance may not be the most ideal distribution matching approach, we encourage researchers in this field to investigate and identify more effective distributional matching methods to ultimately achieve truly optimal VQ strategies. 

This work examines vector quantization (VQ) from a distributional perspective, introducing evaluation criteria that emphasize the importance of aligning feature and code vector distributions. By leveraging the quadratic Wasserstein distance, our approach achieves near-full codebook utilization and significantly reduces quantization error, leading to more stable training and improved image reconstruction performance. These findings underscore the critical role of distributional matching in effective VQ and provide a new direction for designing robust and efficient quantization methods. This insight benefits the broader visual generation community, including applications in image synthesis, data compression, and generative modeling.

From a societal perspective, improvements in VQ contribute to the development of more efficient generative models, potentially reducing the computational costs and energy consumption associated with large-scale training. However, the enhanced visual fidelity enabled by our method, demonstrated particularly on facial datasets like FFHQ, carries inherent risks of misuse, such as the creation of hyper-realistic deepfakes or visual misinformation. Furthermore, because our approach fundamentally relies on data-driven distribution alignment, it may faithfully propagate and entrench demographic biases present in the training datasets. We believe these are well-established challenges in the field of generative AI that require ongoing attention from the research community, including the development of robust forgery detection and careful dataset curation.



% From a societal perspective, improvements in VQ could lead to more efficient generative models, potentially reducing computational costs and energy consumption in large-scale machine learning applications. However, as with any advancements in generative AI, our work may also contribute to both beneficial and potentially concerning applications, such as improved image synthesis techniques that could be used in media production or manipulated content generation. We encourage the research community to consider these broader implications while developing enhanced VQ methodologies.

